\section{Fubini 与 Tonelli：二重积分何时能拆成两次}
\label{sec:12-Real-Analysis-Support/06-Product-Measures-and-Fubini/02}

估计一个随机策略的平均回报时，几乎没人会去算乘积空间上的二重积分。通行做法是先固定状态、对动作的随机性求一次期望，再对状态的分布求一次期望。写代码时也一样：外层循环遍历批次，内层循环遍历样本，谁在外谁在内看着无关紧要。这套操作的合法性来自 Fubini 与 Tonelli 两条定理，而它们并不是无条件成立的——存在一张摆在那里的表格，先按行加得 1，先按列加得 0。

\subsection{有限表格可以随便换序，无穷表格不能}

一张成绩表，按行求和再累加，等于按列求和再累加，因为两边加的是同一个有限集合里的同一批数。加法的交换律和结合律管住了全部情形。

把表格放大到无穷，这条保证就断了。取双下标数组

\[
a_{ij}=
\begin{cases}
1,&j=i,\\
-1,&j=i-1,\\
0,&\text{其他，}
\end{cases}
\qquad i,j\ge1,
\]

即主对角线上放 \(+1\)，紧邻的下方一条对角线上放 \(-1\)。

\begin{center}
\begin{tikzpicture}[>=stealth]
\draw[step=0.9,gray!55,very thin] (0,-4.5) grid (4.5,0);
\foreach \i in {1,...,5}
  \node at ({(\i-1)*0.9+0.45},{-(\i-1)*0.9-0.45}) {\(+1\)};
\foreach \i in {2,...,5}
  \node at ({(\i-2)*0.9+0.45},{-(\i-1)*0.9-0.45}) {\(-1\)};
\node at (5.05,-2.25) {\(\cdots\)};
\node at (2.25,-5.0) {\(\vdots\)};
\node at (5.95,0.42) {行和};
\node at (5.95,-0.45) {\(1\)};
\foreach \i in {2,...,5}
  \node at (5.95,{-(\i-1)*0.9-0.45}) {\(0\)};
\node at (-0.85,-5.45) {列和};
\foreach \j in {1,...,5}
  \node at ({(\j-1)*0.9+0.45},-5.45) {\(0\)};
\node at (5.95,-5.45) {\(1\neq 0\)};
\draw[->,thick] (4.75,-0.45) -- (5.5,-0.45);
\draw[->,thick] (0.45,-4.75) -- (0.45,-5.15);
\end{tikzpicture}
\end{center}

\emph{同一张无穷表格，先按行相加得到 1，先按列相加得到 0；正项与负项各自发散，求和顺序于是成了答案的一部分。}

按行看：第一行只有 \(a_{11}=1\)，往下每一行都是 \(-1\) 与 \(+1\) 相消，行和依次是 \(1,0,0,\dots\)，总和为 1。按列看：第 \(j\) 列有 \(a_{jj}=1\) 与 \(a_{j+1,j}=-1\)，每列和都是 0，总和为 0。两个答案都算得没错，错的是``顺序无所谓''这个默认前提。

裂缝的位置可以指出来：\(\sum_{i,j}\abs{a_{ij}}=\infty\)。正项的总质量是无穷，负项的总质量也是无穷，怎么配对就决定了剩下多少。

\subsection{非负函数不必检查任何条件}

上面那个例子的病根在正负相消。把它去掉，麻烦随之消失。

\begin{theorem}[Tonelli]
设 \((\Omega_1,\mathcal F_1,\mu)\) 与 \((\Omega_2,\mathcal F_2,\nu)\) 为 \(\sigma\)-有限测度空间，\(f:\Omega_1\times\Omega_2\to[0,\infty]\) 对乘积 \(\sigma\)-代数可测。则内层积分作为外层变量的函数可测，且
\[
\int_{\Omega_1\times\Omega_2}\!\!f\,d(\mu\times\nu)
=\int_{\Omega_1}\!\left(\int_{\Omega_2}f(x,y)\,d\nu(y)\right)\!d\mu(x)
=\int_{\Omega_2}\!\left(\int_{\Omega_1}f(x,y)\,d\mu(x)\right)\!d\nu(y).
\]
三者允许同时取 \(+\infty\)，等式仍然成立。
\end{theorem}

证明骨架只有三步。先验证 \(f=\ind_C\) 的情形，此时三个积分分别是 \((\mu\times\nu)(C)\) 与两个截口测度的积分，上一节构造乘积测度时的可数可加性论证正好覆盖这一步。再由线性性推到非负简单函数。最后取一列非负简单函数 \(s_n\uparrow f\)，对内外两层积分各用一次单调收敛定理，把极限搬过积分号。

条件对账：单调收敛只要求非负与递增，这里由 \(f\ge0\) 保证；\(\sigma\)-有限性用在第一步，因为乘积测度的唯一性依赖它；至于结果是否有限，Tonelli 根本不过问——两边同为 \(+\infty\) 也算相等。非负情形下``先算哪一个''从来不必犹豫，原因就在这里。

Tonelli 还附送一条常用的信息：若二重积分有限，则对几乎所有 \(x\)，内层积分 \(\int f(x,y)\,d\nu(y)\) 有限。换句话说，坏的截口至多占一个零测集。

\subsection{带符号时要先给绝对值过一遍 Tonelli}

有正有负的函数不能直接换序，Fubini 的办法是加一道前置检查。

\begin{theorem}[Fubini]
在同样的 \(\sigma\)-有限假设下，若 \(\int\abs{f}\,d(\mu\times\nu)<\infty\)，则两种逐次积分都有定义、彼此相等，并等于二重积分。
\end{theorem}

证明不需要新工具。拆成正部与负部 \(f=f^+-f^-\)，两者非负可测，绝对可积性给出 \(\int f^+<\infty\) 与 \(\int f^-<\infty\)。对二者分别调用 Tonelli，再相减；因为两个数都有限，减法不会撞上 \(\infty-\infty\)。

于是使用流程压缩成一句话：把 \(\abs{f}\) 交给 Tonelli 算一遍，有限就换序，无限就留在二重积分里。回头看那张无穷表格，绝对值之和发散，Fubini 的前提当场失效，两个不同的答案并不构成矛盾。

\subsection{连续版的反例把发散藏在原点附近}

离散反例干净但看着像人为构造。连续情形有一个同样标准的例子，在 \((0,1)^2\) 上取

\[
f(x,y)=\frac{x-y}{(x+y)^3}.
\]

固定 \(y\) 先对 \(x\) 积分。把分子写成 \((x+y)-2y\)，被积函数成为 \((x+y)^{-2}-2y(x+y)^{-3}\)，原函数是 \(-\dfrac{1}{x+y}+\dfrac{y}{(x+y)^2}\)。代入端点，\(x=0\) 处两项恰好抵消，于是

\[
\begin{aligned}
\int_0^1 f(x,y)\,dx
&=-\frac{1}{1+y}+\frac{y}{(1+y)^2}
=-\frac{1}{(1+y)^2},\\[2pt]
\int_0^1\!\left(\int_0^1 f(x,y)\,dx\right)dy
&=\left[\frac{1}{1+y}\right]_0^1=-\frac12 .
\end{aligned}
\]

被积函数在交换两个变量时变号，而积分区域对称，所以另一个顺序给出 \(+1/2\)。

绝对值的积分确实发散，这一点可以粗略但可靠地看出来：在区域 \(0<y<x/2\) 上有 \(x-y>x/2\) 且 \((x+y)^3<8x^3\)，故 \(\abs{f}>1/(16x^2)\)；对 \(y\) 积出长度 \(x/2\) 后剩下 \(1/(32x)\)，再对 \(x\) 在 \((0,1)\) 上积分就是发散的对数积分。发散全部集中在原点附近，而两个方向的逐次积分各自收敛——正是这种``总质量无穷、抵消方式不同''的格局制造了两个答案。

\subsection{期望的嵌套与求和的换序都在用这两条定理}

设 \((X,Y)\) 的联合分布是乘积可测空间上的概率测度 \(P_{X,Y}\)。日常写下的

\[
\E[g(X,Y)]=\E_X\!\bigl[\E[g(X,Y)\given X]\bigr]
\]

在测度语言里就是把二重积分拆成先对 \(Y\) 的条件分布、再对 \(X\) 的边缘分布的逐次积分。条件期望的严格构造见\hyperref[sec:12-Real-Analysis-Support/05-Measure-Integration-and-Conditional-Expectation/05]{``条件期望是对信息 \(\sigma\)-代数的投影''}；这里要强调的是合法性来源：\(g\ge0\) 时 Tonelli 直接放行，\(g\) 带符号时必须先有 \(\E\abs{g(X,Y)}<\infty\)。

独立时的乘积期望是同一件事的特例。若 \(P_{X,Y}=P_X\times P_Y\) 且 \(\E\abs{XY}<\infty\)，则

\[
\begin{aligned}
\E[XY]
&=\iint xy\;d(P_X\times P_Y)
=\int\!\left(\int xy\,dP_Y(y)\right)dP_X(x)\\
&=\left(\int x\,dP_X(x)\right)\!\left(\int y\,dP_Y(y)\right)
=\E[X]\,\E[Y].
\end{aligned}
\]

绝对可积这个前提常被略过，但它不是摆设：重尾场景下 \(\E\abs{XY}\) 可以是无穷，此时等式两边都失去意义。

工程上最容易踩坑的是无穷求和的换序。计算折扣回报的期望时，人们习惯把 \(\E\sum_{t\ge0}\gamma^t r_t\) 写成 \(\sum_{t\ge0}\gamma^t\E[r_t]\)。这一步就是计数测度与概率测度乘积上的 Fubini。奖励非负时 Tonelli 无条件许可；奖励可正可负时，需要 \(\sum_t\gamma^t\E\abs{r_t}<\infty\)，而这正是折扣因子 \(\gamma<1\) 加上奖励有界所提供的保障。折扣一旦取到 1，或者奖励无界，换序的许可证就作废了，\hyperref[sec:16-Control-and-Reinforcement-Learning/02-Reinforcement-Learning/01]{``MDP 把决策问题写成随机过程''}中对折扣的种种要求，一部分动机就在这里。同样的检查也适用于重要性采样权重的期望、无穷级数形式的损失，以及任何``先对样本求和还是先对时间求和''的实现选择。

Tonelli 与 Fubini 不会告诉你怎样把一个难算的二重积分拆成好算的形式，那属于积分技巧。它们只回答一个判定问题：拆完之后答案是否还是原来那个。下一节把这套工具用回随机变量，看独立性怎样落到密度的因式分解上。
